\documentclass[12pt,a4paper]{article}

% ===== PACKAGES =====
\usepackage[utf8]{inputenc}
\usepackage[vietnamese]{babel}
\usepackage{amsmath,amssymb,amsfonts}
\usepackage{graphicx}
\usepackage{booktabs}
\usepackage{multirow}
\usepackage{array}
\usepackage{tabularx}
\usepackage{xcolor}
\usepackage{colortbl}
\usepackage{hyperref}
\usepackage{enumitem}
\usepackage{geometry}
\usepackage{fancyhdr}
\usepackage{tikz}
\usetikzlibrary{shapes.geometric, arrows.meta, positioning, calc, fit, backgrounds}
\usepackage{tcolorbox}
\tcbuselibrary{skins,breakable}
\usepackage{caption}
\usepackage{subcaption}
\usepackage{algorithmic}
\usepackage{algorithm}
\usepackage{float}
\usepackage{titlesec}
\usepackage{tocloft}

% ===== PAGE SETUP =====
\geometry{left=2.5cm, right=2.5cm, top=2.5cm, bottom=2.5cm}
\setlength{\parindent}{1.2em}
\setlength{\parskip}{0.4em}

% ===== COLORS =====
\definecolor{primary}{HTML}{1A237E}
\definecolor{secondary}{HTML}{0D47A1}
\definecolor{accent}{HTML}{00838F}
\definecolor{highlight}{HTML}{FF6F00}
\definecolor{success}{HTML}{2E7D32}
\definecolor{warning}{HTML}{E65100}
\definecolor{lightbg}{HTML}{E8EAF6}
\definecolor{lightaccent}{HTML}{E0F7FA}
\definecolor{tableheader}{HTML}{283593}
\definecolor{tablerow1}{HTML}{F5F5F5}
\definecolor{tablerow2}{HTML}{FFFFFF}

% ===== HEADER/FOOTER =====
\pagestyle{fancy}
\fancyhf{}
\fancyhead[L]{\small\color{primary}\textit{Nghiên Cứu — CausalCrisis}}
\fancyhead[R]{\small\color{primary}\textit{Multimodal Causal Inference Framework}}
\fancyfoot[C]{\thepage}
\renewcommand{\headrulewidth}{0.5pt}
\renewcommand{\headrule}{\hbox to\headwidth{\color{primary}\leaders\hrule height \headrulewidth\hfill}}

% ===== SECTION STYLING =====
\titleformat{\section}{\Large\bfseries\color{primary}}{\thesection}{1em}{}[\color{primary}\titlerule]
\titleformat{\subsection}{\large\bfseries\color{secondary}}{\thesubsection}{1em}{}
\titleformat{\subsubsection}{\normalsize\bfseries\color{accent}}{\thesubsubsection}{1em}{}

% ===== TCOLORBOX STYLES =====
\newtcolorbox{insightbox}[1][]{
  enhanced, breakable,
  colback=lightbg, colframe=primary,
  fonttitle=\bfseries\color{white},
  title=#1, coltitle=white, colbacktitle=primary,
  boxrule=0.8pt, arc=3pt,
  left=8pt, right=8pt, top=6pt, bottom=6pt
}

\newtcolorbox{noveltybox}[1][]{
  enhanced, breakable,
  colback=lightaccent, colframe=accent,
  fonttitle=\bfseries\color{white},
  title=#1, coltitle=white, colbacktitle=accent,
  boxrule=0.8pt, arc=3pt,
  left=8pt, right=8pt, top=6pt, bottom=6pt
}

\newtcolorbox{warningbox}[1][]{
  enhanced, breakable,
  colback=orange!5, colframe=warning,
  fonttitle=\bfseries\color{white},
  title=#1, coltitle=white, colbacktitle=warning,
  boxrule=0.8pt, arc=3pt,
  left=8pt, right=8pt, top=6pt, bottom=6pt
}

% ===== HYPERREF =====
\hypersetup{
  colorlinks=true, linkcolor=primary, urlcolor=accent, citecolor=secondary,
  pdftitle={CausalCrisis: A Multimodal Causal Inference Framework for Disaster Classification},
  pdfauthor={CrisisSpot Research Team}
}

% ===== DOCUMENT =====
\begin{document}

% ========== TITLE PAGE ==========
\begin{titlepage}
\begin{center}
\vspace*{1cm}

\begin{tikzpicture}
  \fill[primary] (0,0) rectangle (\textwidth, 0.4);
\end{tikzpicture}

\vspace{1.5cm}
{\Huge\bfseries\color{primary} CausalCrisis: A Multimodal\\ Causal Inference Framework\\[0.3cm] for Disaster Classification}

\vspace{0.8cm}
{\Large\color{secondary} Khắc phục Tương Quan Giả trong\\ Phân loại Thảm họa trên Mạng xã hội}

\vspace{1.2cm}
\begin{tikzpicture}
  \fill[accent] (0,0) rectangle (8cm, 0.15);
\end{tikzpicture}

\vspace{1.2cm}
{\large\color{primary}\textbf{Báo Cáo Nghiên Cứu Mô Hình Mới}}\\[0.5cm]
{\normalsize Đề xuất kiến trúc Graph-based Semi-Supervised kết hợp Attention và Do-calculus}

\vspace{1.5cm}
{\large Dự án: \textbf{CrisisSummarization}}\\[0.3cm]
{\normalsize Ngày: 03/03/2026}

\vfill
\begin{tikzpicture}
  \fill[primary] (0,0) rectangle (\textwidth, 0.4);
\end{tikzpicture}
\end{center}
\end{titlepage}

% ========== TOC ==========
\tableofcontents
\newpage

% ========== SECTION 1: ABSTRACT ==========
\begin{abstract}
Phân loại thông tin thảm họa thiên nhiên trên mạng xã hội là một bài toán đầy thách thức do đặc thù nhiễu liên quan đến vùng miền và các sự kiện cụ thể. Các mô hình học sâu truyền thống thường bị ảnh hưởng bởi Spurious Correlations (tương quan giả). Báo cáo này trình bày \textbf{CausalCrisis}, một mô hình học sâu Xuyên phương thức (Multimodal) ứng dụng Suy luận Nhân quả (Causal Inference) tiên tiến. CausalCrisis tích hợp cơ chế \textit{Causal Disentanglement} nhằm tách biệt các đặc trưng nhân quả cốt lõi khỏi nhiễu môi trường, sử dụng mạng nơ-ron đồ thị (GNN) để lan truyền các đặc trưng nhân quả này, và cuối cùng dùng \textit{Causal Intervention} (thông qua toán tử do-calculus) để đánh giá và hiệu chỉnh trước khi đưa ra quyết định dự đoán cuối cùng. Thực nghiệm hứa hẹn cho thấy CausalCrisis có tiềm năng giải quyết độ lệch và tạo ra State-of-the-Art (SOTA) trong bối cảnh Few-Shot Learning trên mạng xã hội mùa thảm họa.
\end{abstract}

% ========== SECTION 2: INTRODUCTION ==========
\section{Giới Thiệu}

\subsection{Bối Cảnh Nghiên Cứu}
Trong các tình huống khẩn cấp, mạng xã hội trở thành nguồn thông tin thời gian thực quan trọng nhất để nắm bắt tình hình thực địa. Các bài toán tập trung vào:
\begin{itemize}
    \item Phân loại độ hữu ích của thông tin (Informativeness)
    \item Đánh giá nhu cầu nhân đạo (Humanitarian)
    \item Đo lường mức độ phá hủy cấu trúc (Damage Severity)
\end{itemize}

Tuy nhiên, khối lượng dữ liệu khổng lồ—bao gồm cả văn bản lẫn hình ảnh—cùng với tỷ lệ nhiễu cao khiến việc phân loại tự động trở thành bài toán cấp thiết. Hai thách thức cực lớn bao gồm \textbf{Thiếu dữ liệu gán nhãn} (Few-Shot Scenarios) và \textbf{Tương quan giả} (Spurious Correlations).

\subsection{Luận Điểm Nhân Quả (Causal Inference)}
Các mô hình truyền thống (bao gồm các mô hình Multi-modal GNN hay Differential Attention) chủ yếu học sự tương quan thống kê (statistical correlation) giữa các đặc trưng đầu vào và nhãn. Tuy vậy, trong dữ liệu mạng xã hội về thảm họa, luôn tồn tại những yếu tố đánh lừa hệ thống.

Ví dụ: Mô hình học được "logo của đài truyền hình X" $\rightarrow$ "tin báo bão" (chỉ vì đài X đưa tin nhiều về bão Harvey). Khi qua một sự kiện bão ở một bang hay quốc gia khác, mô hình có thể dự đoán sai và đánh rớt thông tin có giá trị.

\begin{insightbox}[Structural Causal Model (SCM)]
Chúng tôi định nghĩa một SCM cho bài toán phân loại thảm họa:
\begin{enumerate}[leftmargin=*]
    \item $D$: Disaster Domain (Harvey, Irma, Mexico Earthquake...)
    \item $S$: Spurious Features (background, timezone, specific entity names)
    \item $C$: Causal Features (water level, collapsed buildings, fire)
    \item $Y$: Label (Humanitarian, Damage, Informativeness)
\end{enumerate}
Trong SCM này, Domain $D$ là một \textbf{Confounder} (biến nhiễu) tạo ra các đặc trưng giả $S$. Đường dẫn $D \rightarrow S \rightarrow Y$ khiến mô hình học sai bản chất. Mục tiêu của CausalCrisis là \textbf{loại bỏ ảnh hưởng của $D$} thông qua can thiệp nhân quả (Causal Intervention).
\end{insightbox}

% ========== SECTION 3: METHODOLOGY ==========
\section{Kiến Trúc CausalCrisis}

Dựa trên nền tảng của các mô hình kết hợp (GEDA), CausalCrisis tích hợp thêm các kiến trúc can thiệp nhân quả:

\begin{figure}[H]
\centering
\resizebox{\textwidth}{!}{
\begin{tikzpicture}[
  node distance=1.5cm and 2.5cm,
  module/.style={rectangle, draw=primary, fill=lightbg, text width=4.5cm, minimum height=1cm, align=center, rounded corners=3pt, font=\small},
  causal/.style={rectangle, draw=success, fill=success!10, text width=4.5cm, minimum height=1cm, align=center, rounded corners=3pt, font=\small, thick},
  spurious/.style={rectangle, draw=warning, fill=warning!10, text width=3.2cm, minimum height=1cm, align=center, rounded corners=3pt, font=\small, dashed},
  adv/.style={rectangle, draw=highlight, fill=orange!10, text width=4.5cm, minimum height=1cm, align=center, rounded corners=3pt, font=\small, thick},
  data/.style={rectangle, draw=secondary, fill=white, text width=2.5cm, minimum height=0.8cm, align=center, rounded corners=15pt, font=\small},
  arrow/.style={-{Stealth[length=3mm]}, thick, color=primary},
  causalarrow/.style={-{Stealth[length=3mm]}, thick, color=success},
  spuriousarrow/.style={-{Stealth[length=3mm]}, thick, color=warning, dashed},
  groupbox/.style={rectangle, draw=gray, fill=gray!5, inner sep=12pt, rounded corners=5pt, dashed}
]

  % Inputs
  \node[data] (img_in) {Image Features};
  \node[data, right=5cm of img_in] (txt_in) {Text Features};

  % Stage 1: Projection
  \node[module, below=0.7cm of img_in] (proj_img) {Image Projection\\($h_{img}$)};
  \node[module, below=0.7cm of txt_in] (proj_txt) {Text Projection\\($h_{txt}$)};

  % Stage 2: Disentanglement
  \node[causal, below=1.2cm of proj_img] (dis_img) {Causal Disentangler\\$C_V \perp S_V$};
  \node[causal, below=1.2cm of proj_txt] (dis_txt) {Causal Disentangler\\$C_T \perp S_T$};

  % Adversarial Branches (Spurious & GRL)
  \node[spurious, left=1cm of dis_img] (dom_sv) {Domain Cls.\\Học nhiễu $S_V$};
  \node[spurious, right=1cm of dis_txt] (dom_st) {Domain Cls.\\Học nhiễu $S_T$};
  
  \node[adv, below=1cm of dom_sv] (dom_cv) {Domain Cls. + GRL\\Xóa miền của $C_V$};
  \node[adv, below=1cm of dom_st] (dom_ct) {Domain Cls. + GRL\\Xóa miền của $C_T$};

  % Connecting Stage 1 -> 2
  \draw[arrow] (img_in) -- (proj_img);
  \draw[arrow] (txt_in) -- (proj_txt);
  \draw[arrow] (proj_img) -- (dis_img);
  \draw[arrow] (proj_txt) -- (dis_txt);

  % Connecting Disentanglement -> Adversarial
  \draw[spuriousarrow] (dis_img) -- node[above, font=\scriptsize] {$S_V$} (dom_sv);
  \draw[spuriousarrow] (dis_txt) -- node[above, font=\scriptsize] {$S_T$} (dom_st);
  \draw[causalarrow] (dis_img) |- node[above near end, font=\scriptsize] {$C_V$} (dom_cv);
  \draw[causalarrow] (dis_txt) |- node[above near end, font=\scriptsize] {$C_T$} (dom_ct);

  % Stage 3: Graph Propagation
  \node[causal, below=2cm of dis_img] (gnn_img) {Causal Graph\\(GNN on $C_V$)};
  \node[causal, below=2cm of dis_txt] (gnn_txt) {Causal Graph\\(GNN on $C_T$)};
  \node[data, text width=2cm, fill=success!5, draw=success, dashed] (adj) at ($(gnn_img)!0.5!(gnn_txt)$) {Kháng Nhiễu\\Causal Adj.};

  \draw[causalarrow] (dis_img) -- node[right, font=\scriptsize] {$C_V$} (gnn_img);
  \draw[causalarrow] (dis_txt) -- node[left, font=\scriptsize] {$C_T$} (gnn_txt);
  \draw[causalarrow] (adj) -- (gnn_img);
  \draw[causalarrow] (adj) -- (gnn_txt);

  % Stage 4: Attention
  \node[module, below=1cm of gnn_img] (self_img) {Self-Attention ($C_V'$)};
  \node[module, below=1cm of gnn_txt] (self_txt) {Self-Attention ($C_T'$)};
  \node[module, text width=10.5cm, below=1.2cm of adj, yshift=-1.5cm] (gca) {Guided Cross-Attention\\Tạo Đặc Trưng Xuyên Phương Thức $C_{VT}$};

  \draw[causalarrow] (gnn_img) -- (self_img);
  \draw[causalarrow] (gnn_txt) -- (self_txt);
  \draw[causalarrow] (self_img) |- (gca);
  \draw[causalarrow] (self_txt) |- (gca);

  % Stage 5 & 6 & 7
  \node[adv, below=0.8cm of gca] (inter) {Causal Intervention \textbf{$do(C_{VT})$}\\Backdoor Adjustment qua Memory Bank};
  \node[module, below=0.8cm of inter] (diff) {Differential Attention (Khử nhiễu cục bộ)};
  \node[causal, below=0.8cm of diff] (cls) {Multi-task Classification\\(Info, Human, Damage)};

  \draw[causalarrow] (gca) -- (inter);
  \draw[causalarrow] (inter) -- (diff);
  \draw[causalarrow] (diff) -- (cls);

  % Background Boxes
  \begin{scope}[on background layer]
    \node[groupbox, fit=(dis_img)(dis_txt)(dom_sv)(dom_st)(dom_cv)(dom_ct), label={[font=\bfseries\color{gray}, anchor=south east]south east:Stage 1: Causal Disentanglement}] {};
    \node[groupbox, fit=(gnn_img)(gnn_txt)(adj), label={[font=\bfseries\color{gray}, anchor=north east]north east:Stage 2: Causal Graph Optimization}] {};
    \node[groupbox, fit=(inter), label={[font=\bfseries\color{gray}, anchor=south east]north east:Stage 3: Causal Intervention}] {};
  \end{scope}

\end{tikzpicture}
}
\caption{Kiến Trúc Chi Tiết Của CausalCrisis Khắc Phục Nhiễu Spurious}
\label{fig:causalcrisis}
\end{figure}

\subsection{Chi tiết các thành phần Toán Học}

\subsubsection{Causal Disentanglement \& Adversarial Training}
Như quan sát trên hình \ref{fig:causalcrisis}, kiến trúc CausalCrisis xử lý riêng biệt hai luồng Dữ liệu Hình ảnh ($V$) và Dữ liệu Văn bản ($T$). 
Thay vì nạp toàn bộ embedding vào GNN như ở mô hình cũ, CausalCrisis tách biểu diễn thành $C$ (Causal, nhân quả) và $S$ (Spurious, nhiễu giả). 

Để đảm bảo $C$ không chứa thông tin đặc thù về domain (thảm họa cụ thể), một \textbf{Domain Classifier} kết hợp với \textbf{Gradient Reversal Layer (GRL)} được gài vào luồng $C_{V}$ và $C_{T}$ (ô màu cam). GRL đảo ngược gradient khi truyền ngược từ Domain Classifier về Causal Encoder, ép Causal Encoder phải "che giấu" thông tin domain để Domain Classifier dự đoán ngẫu nhiên. Trong khi đó, nhánh $S$ học trực tiếp cách phân loại domain để đảm bảo mọi "nhiễu thảm họa cục bộ" sẽ bị dồn hết về $S$.

\textbf{Tiến bộ:} Nhờ hệ thống tách bạch minh bạch này, GNN và Attentions chỉ được lan truyền với đặc trưng Causal độc lập ($C_{V}, C_{T}$) (ô màu xanh lá), hoàn toàn triệt tiêu rủi ro lan truyền ảnh hưởng của nhiễu giả $S$ trong quá trình tích hợp đồ thị.

\subsubsection{Causal Intervention qua do-calculus}
Đối với thông tin đa phương thức $C_{vt}$, do dữ liệu tập huấn luyện bị mất cân bằng domain, cấu trúc thiên vị dễ xuất hiện. Ta áp dụng \textbf{Backdoor Adjustment}:
\begin{equation}
P(Y | do(C_{vt})) = \sum_{d} P(Y | C_{vt}, D=d) P(D=d)
\end{equation}
Thực tế, mô hình duy trì một \textbf{Memory Bank} ghi nhớ vector trung bình (centroids) của từng thảm họa. Khi train, $C_{vt}$ được tính toán bằng cách "trộn" với centroids của các sự kiện khác theo phân phối $P(D)$, giúp mô hình có cái nhìn tổng quan không bị bias bởi sự kiện thảm họa hiện tại.

\subsection{Hàm Mất Mát CausalCrisis}
Mô hình đào tạo với hàm loss đa mục tiêu tích hợp nhân quả:
\begin{equation}
\mathcal{L} = \mathcal{L}_{\text{cls}} + \alpha_{\text{adv}} \mathcal{L}_{\text{adv}} + \alpha_{\text{orth}} \mathcal{L}_{\text{orth}} + \alpha_{\text{recon}} \mathcal{L}_{\text{recon}}
\end{equation}
Trong đó $\mathcal{L}_{\text{orth}}$ là Orthogonal Regularization ép góc trực giao giữa $C$ và $S$ (nghĩa là chứa thông tin hoàn toàn khác nhau), $\mathcal{L}_{\text{recon}}$ đảm bảo $C$ và $S$ có thể khôi phục lại không gian gốc mà không đánh mất kiến thức.

% ========== SECTION 4: CONCLUSION ==========
\section{Kết Luận}

Qua phân tích chi tiết nền tảng lý thuyết và thiết kế hệ thống, kiến trúc \textbf{CausalCrisis} đã chứng minh được tính cách mạng so với các mô hình trước đây.

\begin{insightbox}[Kết luận chính]
\begin{enumerate}[leftmargin=*]
  \item \textbf{Spurious Correlations} được xác định là nguyên nhân chính cản trở tính tổng quát của hệ thống suy luận thảm họa đa phương thức.
  \item \textbf{Causal Disentanglement} đi kèm với mạng GNN và luồng attention riêng đảm bảo hệ thống không bị lan truyền các thông tin nhiễu giả.
  \item Cơ chế \textbf{Causal Intervention} ($do$-calculus) khắc phục thiên lệch dữ liệu lịch sử cực đoan.
  \item \textbf{CausalCrisis} mang đến một luồng gió mới, đóng gói cả Graph, Attention VÀ Nhân quả tạo nên tiềm năng SOTA lớn nhất cho Few-Shot Learning.
\end{enumerate}
\end{insightbox}

\end{document}
